\documentclass[12pt]{article}
\usepackage[utf8]{inputenc}
\usepackage[T1]{fontenc}
\usepackage{hyperref}
\title{Havasupai-Hualapai-Yavapai Language Resources for K-12 Curriculum Development}
\author{Compiled by Claude Sonnet 5.5 (claude-sonnet-5-5) for Andr\'as Kornai}
\date{}
\begin{document}
\maketitle

\section{Introduction}
Havasupai-Hualapai-Yavapai (yuf) is a Cochimi-Yuman language of Arizona. The Havasupai and Hualapai dialects use different orthographies. The searches found more on Hualapai than on Havasupai or Yavapai.

\section{Primary Educational Resources}
Lucille Watahomigie worked on bilingual curriculum at Peach Springs School. A more recent source reports a bilingual programme there plus immersion camps. Related material: \url{https://www8.nau.edu/jar/SIL/SchoolsDev.pdf}. No specific textbook was identified.

\section{Linguistic and Reference Materials}
\emph{A Dictionary of the Hualapai Language} (2003): \url{https://knowledgebase.arts.ubc.ca/?p=3583}, \url{https://wiki.ubc.ca/Documentation:RelLex/A_Dictionary_of_the_Hualapai_Language}. Overviews: \url{https://en.wikipedia.org/wiki/Havasupai%E2%80%93Hualapai_language}, \url{https://en.wikipedia.org/wiki/Yavapai_language} (which lists a Yavapai phonology and dictionary thesis).

\section{Religious and Community Resources}
None were found. A community essay on Havasupai: \url{https://thelanguagecloset.com/2025/01/25/havasupai/}.

\section{Digital and Audio Resources}
Endangered Languages Project entry: \url{https://endangeredlanguages.com/node/94360}.

\section{Curriculum Development Context}
Hualapai has the most documentation and school involvement. Havasupai and Yavapai materials would need separate searches, for example for the Supai school and the Yavapai-Prescott and Fort McDowell programmes.

\end{document}